\documentclass[a4paper,oneside,12pt]{book} % chcete-li si VÚ vytisknout, změňte "oneside" na "twoside"
%% === nezbytné balíčky:
%\usepackage[T1]{fontenc}
\usepackage[IL2]{fontenc}   % písmo pro UTF-8
\usepackage[utf8]{inputenc}    % vstupní znaková sada UTF-8
%\usepackage[cp1250]{inputenc}   % vstupní znaková sada: Windows 1250

\usepackage[czech]{babel} % česky psaná práce, typografická pravidla     
%\usepackage{encxvlna} % postará se o spojky a předložky, které dle českých pravidel nesmí být na konci řádku. Dokumentace: http://texdoc.net/texmf-dist/doc/generic/encxvlna/encxvlna.pdf
\usepackage[a4paper, hmarginratio=2:2]{geometry} % využití celé A4 stránky a nastavení okrajů. Chcete-li si VÚ vytisknout, změňte: hmarginratio=3:2
\usepackage{pdfpages} % pokud nemáte formulář "Zadání VÚ" naskenovaný jako PDF, tak ZAKOMENTUJTE

%% === balíčky, které se mohou hodit:
\usepackage[hidelinks]{hyperref} % v PDF budou klikací odkazy ("hidelinks" je nebude rámovat)
\usepackage{graphicx} % balíček pro vkládání __rastrových__ grafických souborů (PNG apod.)
%\usepackage{epsfig} % balíčky pro vkládání grafických souborů typu EPS
%\usepackage{float} % rozšířené možnosti umístění obrázků

%\usepackage{caption} % pro popisky obrázků, tabulek atd.

\usepackage{tabularx} % rozšířené možnosti tabulek
%\usepackage{tabu} % jiný balík pro rozšířené možnosti tabulek

\usepackage{listings}  % balíček vhodný pro ukázky kódu 
\usepackage{amsmath} % balíček pro pokročilou matematickou sazbu
%\usepackage{color} % pro možnost barevného textu
%\usepackage{fancybox} % umožňuje pokročilé rámečkování

%\usepackage{index} % nutno použít v případě tvorby rejstříku balíčkem makeindex
%\newindex{default}{idx}{ind}{Rejstřík} % zavádí rejstřík v případě použití balíku index


\topmargin=-15mm      % horní okraj trochu menší
\textwidth=150mm      % šířka textu na stránce
\textheight=240mm     % "výška" textu na stránce

\frenchspacing % za větou bude mezislovní mezera (v anglických textech je mezera za větou delší)
\widowpenalty=1000 % "síla" zákazu vdov (= jeden řádek odstavce na konci stránky)
\clubpenalty=1000 % "síla" zákazu sirotků (= jeden řádek/slovo odstavce samostatně na začátku stránky)
\brokenpenalty=1000 % "síla" zákazu zlomu stránky za řádkem, který má na konci rozdělené slovo

\pagenumbering{arabic} % číslování stránek arabskými číslicemi
\pagestyle{plain}      % stránky číslované dole uprostřed

\parindent=0pt % odsazení 1. řádku odstavce
\parskip=7pt   % mezera mezi odstavci

\newcommand{\ti}{\textit} % zkrácený příkaz pro kurzívu
\newcommand{\tb}{\textbf} % zkrácený příkaz pro tučné písmo


%% --- zde jsou zavedeny některé "konstanty" - některé musíte změnit! --- %%
\newcommand{\cvut}{České vysoké učení technické v~Praze}
\newcommand{\fjfi}{Fakulta jaderná a fyzikálně inženýrská}
\newcommand{\ksi}{Katedra softwarového inženýrství}
\newcommand{\program}{Aplikace informatiky v~přírodních vědách} % změňte, pokud máte jiný stud. program
\newcommand{\specializace}{--} % změňte, pokud máte specializaci stud. programu

\newcommand{\druh}{Výzkumný úkol} % druh práce
\newcommand{\woman}{} % pokud jste ŽENA, ZMĚŇTE na: ...{\woman}{a} (je to do Prohlášení)

\newcommand{\logoCVUT}{\includegraphics{symbol_cvut_konturova_verze_cb.pdf}} % logo ČVUT -- podle grafického manuálu ČVUT platného od prosince 2016. Pokud nevyhovuje PDF-verze, tak použijte jinou variantu loga: https://www.cvut.cz/logo-a-graficky-manual -> "Symbol a logo ČVUT v Praze"). Pokud chcete logo úplně vynechat, zadejte místo "\includegraphics{...}" text "\vspace{35mm}"

% přesně podle formuláře "Zadání VÚ" VYPLŇTE:
\newcommand{\nazevcz}{Český název VÚ}    % český název práce (přesně podle zadání!)
\newcommand{\autor}{Jméno Příjmení}   % vyplňte své jméno a příjmení (s akademickým titulem, máte-li jej)
\newcommand{\vedouci}{JménoV PříjmeníV} % vyplňte jméno a příjmení vedoucího práce, včetně titulů, např.: Doc. Ing. Ivo Malý, Ph.D.
\newcommand{\pracovisteVed}{\ksi, \fjfi, \cvut} % ZMĚŇTE, pokud vedoucí Vaší práce není z KSI
\newcommand{\konzultant}{--} % POKUD MÁTE určeného konzultanta, NAPIŠTE jeho jméno a příjmení
\newcommand{\pracovisteKonz}{--} % POKUD MÁTE konzultanta, NAPIŠTE jeho pracoviště

% podle skutečnosti VYPLŇTE:
\newcommand{\rok}{2023}  % rok odevzdání práce (jen rok odevzdání, nikoli celý akademický rok!)
\newcommand{\kde}{Praze}

\newcommand{\klicova}{Klíčová slova}   % zde NAPIŠTE česky max. 5 klíčových slov
\newcommand{\abstrCZ}{Popis VÚ česky}    % zde NAPIŠTE abstrakt v češtině (cca 7 vět, min. 80 slov)

\newcommand{\prohlaseni}{Prohlašuji, že jsem svůj výzkumný úkol vypracoval\woman{} samostatně a použil\woman{} jsem pouze podklady (literaturu, projekty, SW atd.) uvedené v přiloženém seznamu.} % text prohlášení můžete mírně upravit :-)

\newcommand{\podekovani}{Děkuji ... za ...} % NAPIŠTE poděkování, např. svému vedoucímu:
%"Děkuji Ing. Eleonoře Krtečkové, Ph.D. za vedení mého výzkumného úkolu a za podnětné návrhy, které jej obohatily."
% NEBO:
% Děkuji doc. Pafnutijovi Snědldítětikaši, Ph.D. za neocenitelné rady a pomoc při tvorbě výzkumného úkolu.


\begin{document}
%%%%%%%%%%%% TITULNÍ STRANA -- následujících cca 30 řádků se generuje AUTOMATICKY (není třeba nic měnit) %%%%%%%%%%%%
\thispagestyle{empty}

\begin{center}
    {\Large \textsc{\cvut}\\[1.5ex] \textsc{\fjfi}}\\
    \vspace{10mm}

    \begin{tabular}{ll}
		\multicolumn{2}{c}{\tb{\ksi}} \\[8pt]   
		\tb{Studijní program:} & \tb{\program}\\
		\tb{Specializace:} & \tb{\specializace}\\
    \end{tabular}

   \vspace{13mm} \logoCVUT \vspace{15mm} 

   {\huge \tb{\nazevcz}\par}
   \vspace{5mm}   
   
   \vspace{15mm}
   {\Large \MakeUppercase{\druh}}

   \vfill
   {\large
    \begin{tabular}{ll}
    Vypracoval: & \autor\\
    Vedoucí práce: & \vedouci\\
    Rok: & \rok
    \end{tabular}
   }
\end{center}

%%%%%%%%%%%% ZADÁNÍ PRÁCE %%%%%%%%%%%%
% Zadání (podepsané!) musíte NASKENOVAT. Ideálně jako jednostránkové PDF (soubor "zadaniVU.pdf"). 
%\newpage  % SEM NESAHEJTE!
\thispagestyle{empty} % SEM NESAHEJTE!

% --- varianta A: zadání naskenované jako 2stránkové PDF:
\includepdf[pages={1}]{zadaniVU.pdf} % zadání naskenované jako PDF
%% --- varianta B: zadání naskenované jako obrázek (JPG):
%\begin{center}
%     \includegraphics[width=1\textwidth]{zadani1.jpg}
%\end{center}
%\newpage ~  % prázdná stránka za zadáním -- oboustranný tisk; SEM NESAHEJTE!
%\thispagestyle{empty} % SEM NESAHEJTE!


%%%%%%%%%%%% Prohlášení -- SEM NESAHEJTE! Generuje se automaticky z výše nastavených maker \kde{} a \prohlaseni{}. %%%%%%%%%%%%
\newpage % SEM NESAHEJTE!
\thispagestyle{empty}  % SEM NESAHEJTE!

~ % SEM NESAHEJTE!
\vfill % prázdné místo. SEM NESAHEJTE!

\tb{Prohlášení} % SEM NESAHEJTE!

\vspace{1em} % vertikální mezera. SEM NESAHEJTE!
\prohlaseni

\vspace{2em}  % SEM NESAHEJTE!
\hspace{-0.5em}\begin{tabularx}{\textwidth}{X c}  % SEM NESAHEJTE!
V \kde\ dne .................... &........................................ \\	% SEM NESAHEJTE!
	& \autor
\end{tabularx}	% SEM NESAHEJTE!


%%%%%%%%%%%% Poděkování -- tuto stránku můžete celou odstranit %%%%%%%%%%%%
\newpage
\thispagestyle{empty}

~
\vfill % prázdné místo

\tb{Poděkování}

\vspace{1em} % vertikální mezera
\podekovani
\begin{flushright}
\autor
\end{flushright}  % <-- konec stránky s poděkováním


%%%%%%%%%%%% ABSTRAKT atp. Je generován AUTOMATICKY podle maker nastavených na začátku souboru) %%%%%%%%%%%% 
\newpage   % SEM NESAHEJTE!
\thispagestyle{empty}   % SEM NESAHEJTE!

% příprava:    (na následujících 8 řádků NESAHEJTE!)
\newbox\odstavecbox
\newlength\vyskaodstavce
\newcommand\odstavec[2]{%
    \setbox\odstavecbox=\hbox{%
         \parbox[t]{#1}{#2\vrule width 0pt depth 4pt}}%
    \global\vyskaodstavce=\dp\odstavecbox
    \box\odstavecbox}
\newcommand{\delka}{120mm} % šířka textů ve 2. sloupci tabulky

% použití přípravy:    % dovnitř "tabular" vůbec NESAHEJTE!
\begin{tabular}{ll}
  {\em Název práce:} & ~ \\
  \multicolumn{2}{l}{\odstavec{\textwidth}{\bf \nazevcz}} \\[1em]
  {\em Autor:} & \autor \\[1em]
  {\em Studijní program:} & \program \\
  {\em Specializace:} & \specializace\\
  {\em Druh práce:} & \druh \\[1em]
  {\em Vedoucí práce:} & \odstavec{\delka}{\vedouci\\ \pracovisteVed} \\
  {\em Konzultant:} & -- %\odstavec{\delka}{\konzultant \\ \pracovisteKonz}  % VYMAŽTE text "-- %" v případě, že jste neměli konzultanta
 \\[1em]  
  \multicolumn{2}{l}{\odstavec{\textwidth}{{\em Abstrakt:} ~ \abstrCZ  }} \\[1em]
  {\em Klíčová slova:} & \odstavec{\delka}{\klicova}
\end{tabular}



%%%%%%%%%%%% Obsah VÚ ... je generován AUTOMATICKY %%%%%%%%%%%%
\parskip=0pt
\tableofcontents % SEM NESAHEJTE!
\parskip=7pt
%\newpage % SEM NESAHEJTE!


%--------------------------------------------------------
%|         Zde začíná SAMOTNÁ PRÁCE (text)              |
%--------------------------------------------------------

\chapter*{Úvod} % SEM NESAHEJTE!
\addcontentsline{toc}{chapter}{Úvod} % SEM NESAHEJTE!
%
Zde napište text úvodu (1-3 strany, nerozdělujte na podkapitoly) nebo jej vložte ze samostatného souboru: např. příkazem \texttt{\textbackslash input\{uvod.tex\}}.

Tato šablona je určena {\color{red}pro elektronické odevzdání VÚ} (od roku 2023). Pokud byste si VÚ chtěli vytisknout, tak změňte řádky~1~a~10 ve zdrojovém kódu této šablony ($\Rightarrow$ budou správně nastavené \uv{prázdné zadní stránky}).
%
%\input{uvod.tex}



\chapter{Název první kapitoly}
%
Tady začněte psát text první kapitoly práce nebo jej vložte ze samostatného souboru, např. příkazem \texttt{\textbackslash input\{kapitola1.tex\}}.
%
%\input{kapitola1.tex} % text vkládán ze souboru. Pokud je v souboru uveden i příkaz \chapter{...}, tak ho o 4 řádky výše vymažte!


\chapter*{Závěr}
\addcontentsline{toc}{chapter}{Závěr} % SEM NESAHEJTE!
%
Zde napište text závěru (1-3 strany, nerozdělujte na podkapitoly) nebo jej vložte ze samostatného souboru: např. příkazem \texttt{\textbackslash input\{zaver.tex\}}.
%
%\input{zaver.tex}


%%%%%%%%%%%% SEZNAM POUŽITÝCH ZDROJŮ (LITERATURA) %%%%%%%%%%%%
\addcontentsline{toc}{chapter}{Literatura} % SEM NESAHEJTE!
\begin{thebibliography}{99}   
	\bibitem{odkaz} Autor. \ti{Název knihy}. Město. Nakladatelství. Rok.  
\end{thebibliography}
% formát: ČSN ISO 690. Můžete si to vygenerovat na http://www.citacepro.com (přihlaste se přes odkaz "ČVUT"), umí to vygenerovat TeX
% řazení: abecedně podle autora (resp. provního slova, není-li znám autor)


%%%%%%%%%%%% PŘÍLOHY PRÁCE %%%%%%%%%%%%
\newpage % SEM NESAHEJTE!
\addcontentsline{toc}{chapter}{Přílohy} % SEM NESAHEJTE!
\appendix % SEM NESAHEJTE!


%%%%%%%%%%%% Příloha A (tj. 1. kapitola v rámci příloh) %%%%%%%%%%%%
\chapter{Název přílohy} % SEM NESAHEJTE!
%
Zde napište text první přílohy nebo jej vložte, např.: \texttt{\textbackslash input\{priloha\_A.tex\}}.
%
%\input{priloha_A.tex} % text vkládán ze souboru, kde je i příkaz \chapter{...}


\end{document} % SEM NESAHEJTE! Konec.
